% Reeksen. Hammershøi: deur na deur, elke opening kleiner, en toch blijft alles binnen één ruimte.
% Motief: een vierkant met zijde 1, opgevuld met stukken van 1/2, 1/4, 1/8, ...: de meetkundige
% reeks 1/2 + 1/4 + 1/8 + ... = 1. Elk stuk is de helft van wat overblijft, afwisselend verticaal
% en horizontaal gesneden (tien stukken getekend; de rest is kleiner dan een millimeter).
\def\omslagbestand{omslag-deuren.jpg}
\def\omslagbron{Vilhelm Hammershøi, \textit{Witte deuren, open deuren}, 1905 (publiek domein)}
\def\omslagrand{10.45 + 0.45*\x/21}
\def\omslagmotief{%
  \randlijn
  \begin{scope}[shift={(14.6,4.1)},scale=5.4]
    % stuk k: [x0,x1] x [y0,y1]; het overblijvende deel wordt afwisselend links en boven gehalveerd
    \def\xa{0}\def\ya{0}\def\xb{1}\def\yb{1}%
    \foreach \k in {1,...,10}{%
      \pgfmathsetmacro{\dek}{max(.2,1-.16*(\k-1))}%
      \ifodd\k
        \pgfmathsetmacro{\xm}{(\xa+\xb)/2}%
        \fill[fel,opacity=\dek] (\xa,\ya) rectangle (\xm,\yb);
        \draw[white,line width=.9pt] (\xm,\ya) -- (\xm,\yb);
        \xdef\xa{\xm}%
      \else
        \pgfmathsetmacro{\ym}{(\ya+\yb)/2}%
        \fill[fel,opacity=\dek] (\xa,\ym) rectangle (\xb,\yb);
        \draw[white,line width=.9pt] (\xa,\ym) -- (\xb,\ym);
        \xdef\yb{\ym}%
      \fi}
    \draw[omslaglijn=1.4pt] (0,0) rectangle (1,1);
    \node[text=inkt,font=\LARGE] at (.25,.5) {$\displaystyle\frac12$};
    \node[text=inkt,font=\Large] at (.75,.75) {$\displaystyle\frac14$};
    \node[text=inkt,font=\normalsize] at (.625,.25) {$\displaystyle\frac18$};
  \end{scope}
}
